\def\title{\large{چکیده}}
\def\keyword{واژه‌های کلیدی: }
\def\abstractfa{ 
	
	امروزه شبکه‌های تطبیقی توزیع‌شده به دلیل مقیاس‌پذیری بالا، مقاومت در برابر خرابی‌ها و غیرمتمرکز بودن بسیار مورد توجه محققان در زمینه یادگیری توزیع‌شده قرار گرفته‌اند. فرآیند یادگیری الگوریتم‌های توزیع‌شده باید دارای ساختاری مشخص بوده و مبتنی بر یک مدل محاسباتی موازی انجام شود تا ارتباط بین نرم‌افزار موازی و سخت‌افزار موازی به راحتی انجام گردد. مدل محاسباتی موازی \lr{BSP} به راحتی قابل اعمال بر روی این شبکه‌ها است؛ اما در یک محیط اجرایی ناهمگن در صورت اعمال همگام‌سازی مانعی در \lr{BSP}، مساله کندکننده‌ها رخ می‌دهد و میانگین زمان انتظار گره‌ها افزایش می‌یابد. در این رساله به منظور مدیریت محیط ناهمگن، در مرحله اول با استفاده از یک روش خو‌شه‌بندی محتوایی-ساختاری گره‌های شبکه منطقه‌بندی می‌شوند. در مرحله دوم یک مدل محاسباتی موازی توسعه‌یافته به نام \lr{BSP} فدرالی یا \lr{FBSP} ارائه شده است. مدل محاسباتی پیشنهادی همانند \lr{BSP} از سه بخش مهم محاسبه، ارتباط و همگام‌سازی مانعی تشکیل شده است با این تفاوت که در \lr{FBSP}، دو نوع ابرگام منطقه‌ای و ابرگام سراسری وجود دارد. همچنین، ارتباط گره‌ها و همگام‌سازی مانعی مدل‌های تخمین‌زده شده توسط گره‌ها در دو سطح منطقه‌ای و سراسری انجام می‌شود. ارتباط دو سطحی سبب کاهش هزینه‌های ارتباطی و همگام‌سازی مانعی دو سطحی سبب مقابله با مساله کندکننده‌ها و کاهش میانگین زمان انتظار گره‌ها خواهد شد. در مرحله بعد، یک رویکرد یادگیری توزیع‌شده مبتنی بر مدل \lr{FBSP} برای حل مسائل بهینه‌سازی ارائه شد. منظور مقاوم‌سازی این رویکرد در برابر انواع نویزها و سازگاری با محیط‌های ایستا و غیرایستا، الگوریتم یادگیری توزیع‌شده انتشاری به نام \lr{RDSGD} مبتنی بر گرادیان کاهشی تصادفی معرفی شد. در نهایت، ابزار نرم‌افزاری مستقلی به نام \lr{FBSPonMPI} برای توسعه برنامه‌های توزیع‌شده در بستر شبکه‌های تطبیقی به زبان \lr{C} و پایتون ارائه شده است که سریع و مستقل از سخت‌افزار است. تحلیل همگرایی و کارایی رویکرد پیشنهادی در حالت میانگین و مربع میانگین و بررسی تاثیر توپولوژی تصادفی در شبکه و خطا کمی‌سازی انجام شده است. جهت تصدیق نتایج تحلیلی بدست آمده، رویکرد پیشنهادی با استفاده از تعدادی معیار ارزیابی با سایر رویکرد‌ها مقایسه شده است. رویکرد مبتنی بر مدل \lr{FBSP} بر روی مساله بهینه‌سازی شناسایی سیستم و پیش‌بینی بازار سهام مورد ارزیابی قرار گرفت. رویکرد پیشنهادی در حضور نویزهای مختلف و در محیط‌های ایستا و غیرایستا نسبت به سایر روش‌ها مقاوم‌تر و سازگارتر است. نتایج بدست آمده نشان می‌دهند که رویکرد پیشنهادی موفق به کاهش $36.9\%$ در میانگین زمان انتظار و کاهش $58\%$ در میانگین تعداد ارتباطات با حفظ نرخ همگرایی و دقت یادگیری شده است.
}
\def\keywordfa{شبکه‌های تطبیقی‌ توزیع‌شده، یادگیری توزیع‌شده، مدل‌ محاسباتی موازی، مدل محاسباتی همگام مانعی فدرالی، ابرگام منطقه‌ای، ابرگام سراسری}

\begin{center}
\title
\end{center}
\abstractfa

\textbf{واژه‌های کلیدی:}
\keywordfa